\section{Free Maxwell fields}\label{sec:maxwell}

The limiting fields of Section~\ref{sec:infinite-volume} are described by their correlations, meaning the expectations of products of local polynomial measurements.
For a short-distance comparison we use the free Maxwell field: the linear, non-interacting theory of a curvature two-form. It has one component for each basis direction in the Lie algebra of Definition~\ref{def:gauge}; the index~$c$ below records that basis direction (often called colour). Its local observables are polynomials in the curvature and its derivatives. The two linear Maxwell equations say, respectively, that the curvature has zero cyclic derivative and zero divergence. In coordinates, the first says the signed sum around each three-index cycle vanishes; the second says each column has zero divergence.

Subsection~\ref{subsec:mx-curvature} defines the Taylor coefficients of a Maxwell curvature and computes their free correlations by differentiating the curvature two-point function.
Subsection~\ref{subsec:mx-positive} starts with positive pairings of reflected linear observations. To see the finite matching rule, use independent standard complex Gaussian coordinates~$z_a$, each with density~$\pi^{-1}e^{-|z_a|^2}$: they satisfy~$\mathbb E[z_a\overline z_b]=\delta_{ab}$ and~$\mathbb E[z_az_b]=0$. Thus a product has nonzero expectation only when each unbarred factor is paired with a barred factor of the same coordinate. The field correlations below are Wick ordered: pairings within one insertion are removed, leaving the matching rule between distinct insertions. This representation gives polynomial positivity, identifies coefficients that represent the same field, and supplies finitely many detecting tests at each degree.
Subsection~\ref{subsec:mx-comparison} states the additional same-family lattice estimate needed to compare free correlations with the already chosen continuum correlations.
Subsection~\ref{subsec:mx-bases} makes the comparison maps agree across degrees and shows they are isomorphisms.
Subsection~\ref{subsec:mx-smooth} upgrades convergence of tested correlations to convergence of all derivatives in the insertion positions.
The radii and field maps are chosen together. The further uniformity needed for a prescribed radius sequence is left to the later trace argument.


\subsection{Curvature correlations}\label{subsec:mx-curvature}

Use the gauge group of Definition~\ref{def:gauge}, the classical polynomial spaces~$V_d$ of Definition~\ref{def:uv-labels}, and their norms from Definition~\ref{def:iv-coefficients}.
The inner product on the Lie algebra and a corresponding orthonormal basis~$E^1,\ldots,E^{\dim\g}$ are fixed by
\begin{equation}\label{eq:mx-colors}
 \langle X,Y\rangle_{\g}=-\frac{1}{m}\operatorname{Re}\tr(XY),
 \qquad \langle E^c,E^{c'}\rangle_{\g}=\delta_{cc'}.
\end{equation}
Here~$m$ is the defining matrix size, as in Definition~\ref{def:gauge}. Thus $c$ records which Lie-algebra component of the curvature is being measured; the normalisation fixes the trace convention in every later two-point function.\sourcefootnote{\nolinkurl{YM/Gauge/Colors.lean}, lines 43--83: physical colour coordinates and normalised trace}

\begin{definition}[Maxwell Taylor coefficients]\label{def:mx-jets}
The indices~$a,b,k$ range over the four coordinate directions.
A Maxwell field is a smooth array~$f=(f_{ab})$ of~$\g$-valued functions. Antisymmetry makes it a two-form: each component is attached to an oriented coordinate area, and reversing the area changes its sign. The cyclic equation is the zero exterior derivative condition, and the last equation is zero divergence. It satisfies
\begin{equation}\label{eq:mx-equations}
 f_{ab}=-f_{ba},\qquad
 \partial_af_{bk}+\partial_bf_{ka}+\partial_kf_{ab}=0,\qquad
 \sum_a\partial_af_{ab}=0.
\end{equation}
For~$j\in\N$, let~$\mathcal H_j$ be the finite-dimensional real space of homogeneous polynomial solutions of degree~$j$; these are the possible degree-$j$ Taylor terms of a Maxwell curvature near the origin.
A homogeneous polynomial~$Q_j$ is identified with its symmetric multilinear coefficient by evaluation on repeated arguments: its value at~$z$ is obtained by putting the same vector~$z$ in every Taylor slot.
Our Taylor convention is
\begin{equation*}
 Q_j(z)=\frac{1}{j!}D^jf(0)[z,\ldots,z].
\end{equation*}
Thus a Taylor tensor here includes the factorial denominator, so the sequence $(Q_j)$ is exactly the formal Taylor expansion of the curvature.
A sequence~$Q=(Q_j)_{j\ge0}$ with~$Q_j\in\mathcal H_j$ is a formal Maxwell Taylor series; convergence of this series is not required.
For a one-form~$A$, use~$z\cdot A=\sum_a z^aA_a$ and~$(dA)_{ab}=\partial_aA_b-\partial_bA_a$. The radial potential is a polynomial one-form reconstructed from the curvature, with~$dA=Q$. The condition~$z\cdot A(z)=0$ fixes the radial gauge; its explicit formula is
\begin{equation}\label{eq:mx-radial}
 (\kappa_jQ)_a(z)=\frac{1}{j+2}\sum_b z^bQ_{ba}(z),
 \qquad z\cdot\kappa_jQ(z)=0,\qquad d(\kappa_jQ)=Q.
\end{equation}
For example, when~$Q_0$ is a constant antisymmetric matrix, the formula gives~$(\kappa_0Q_0)_a(z)=\frac12\sum_bz^b(Q_0)_{ba}$, a linear one-form whose curl is~$Q_0$.
\end{definition}

The last identity follows by differentiating the formula and using the cyclic equation in~\eqref{eq:mx-equations}. Euler’s identity, $z\cdot\nabla Q_j=jQ_j$, supplies the factor~$j+2$ in the denominator.
Thus the coefficients may also be read in the radial coordinates of~\eqref{eq:uv-radial}.\sourcefootnote{\nolinkurl{YM/Maxwell/Shell.lean}: the constrained Taylor space and inverse radial curl}

For~$z\in\R^4\setminus\{0\}$, the scalar Green function~$\Gamma$ is the fundamental solution of the four-dimensional Laplacian. Its second derivatives are the two-point covariance: they give the free expectation of two linear curvature measurements at separation~$z$.
The indices~$a,b,e,f$ specify coordinate directions; $c,c'$ are the colour indices of~\eqref{eq:mx-colors}.
Set
\begin{equation}\label{eq:mx-kernel}
\begin{aligned}
 \Gamma(z)&=\frac{1}{4\pi^2|z|^2},\qquad z\ne0,\\
 C^{cc'}_{ab,ef}(z)&=\delta_{cc'}\bigl(
 \delta_{be}\partial_a\partial_f\Gamma
 -\delta_{bf}\partial_a\partial_e\Gamma
 -\delta_{ae}\partial_b\partial_f\Gamma
 +\delta_{af}\partial_b\partial_e\Gamma\bigr)(z).
\end{aligned}
\end{equation}
The signs include differentiation at the second endpoint of a covariance depending on~$x-y$.
A linear probe of order~$j$ is a real linear functional on~$\mathcal H_j$; it measures a specified linear combination of the degree-$j$ Taylor coefficients.
For probes~$p\in\mathcal H_j^*$ and~$q\in\mathcal H_\ell^*$, the dual spaces of these real tensor spaces, define
\begin{equation}\label{eq:mx-leg}
 K_{p,q}(z)=\frac{(-1)^\ell}{j!\,\ell!}
 (p\otimes q)\bigl[D^{j+\ell}C(z)\bigr].
\end{equation}
The first~$j$ derivative slots belong to~$p$ and the remaining~$\ell$ slots to~$q$. Thus~$p$ and~$q$ select one Taylor coefficient at each of the two points; the factorials match the Taylor convention in~\eqref{eq:mx-radial}.
Colour and form indices are contracted by the two functionals themselves; no extra antisymmetry factor is inserted.
One may extend these functionals to all tensors to evaluate the formula. The next lemma shows that the equations force the differentiated covariance into the Maxwell solution spaces, so the chosen extensions cannot affect the answer.

\begin{lemma}[The contraction kernel]\label{lem:mx-kernel}
For~$z\ne0$, the kernel~\eqref{eq:mx-kernel} satisfies~\eqref{eq:mx-equations} in both groups of tensor indices.
Consequently~\eqref{eq:mx-leg} is well defined on~$\mathcal H_j^*\times\mathcal H_\ell^*$, is real analytic, and satisfies
\begin{equation}\label{eq:mx-leg-laws}
 K_{p,q}(z)=K_{q,p}(-z),\qquad
 K_{p,q}(\lambda z)=\lambda^{-j-\ell-4}K_{p,q}(z)\quad(\lambda>0).
\end{equation}
\end{lemma}

\begin{proof}
Here the curl is the antisymmetric derivative~$d$ defined above; the cyclic equation says its next exterior derivative vanishes, while harmonicity of~$\Gamma$ gives zero divergence away from the origin.
The two differential equations follow by differentiating the explicit covariance: antisymmetry and commuting derivatives give the cyclic equation, while harmonicity of~$\Gamma$ gives zero divergence away from the origin. Endpoint exchange gives the equations in the other tensor indices.

For the divergence calculation, write~$\Delta=\sum_a\partial_a^2$ for the Laplacian and suppress the colour factor:
\begin{equation*}
 \sum_a\partial_a C_{ab,ef}
 =\delta_{be}\partial_f\Delta\Gamma-\delta_{bf}\partial_e\Delta\Gamma=0
 \qquad(z\ne0).
\end{equation*}
Antisymmetry is explicit in~\eqref{eq:mx-kernel}.
Exchanging the endpoints proves the same equations in the other group of indices.
Every Taylor tensor to which an extended functional is applied therefore lies in~$\mathcal H_j$ or~$\mathcal H_\ell$.
An extension vanishing on that space contributes zero, proving independence.
The exchange and scaling identities follow by differentiating the even, homogeneous function~$\Gamma$; analyticity follows from its explicit formula.\sourcefootnote{\nolinkurl{YM/Maxwell/Curls.lean}: equations in both groups of tensor indices; \nolinkurl{YM/Maxwell/Leg.lean}, lines 141--176: independence of the extensions}
\end{proof}

\begin{definition}[Free polynomial fields]\label{def:mx-fields}
Choose coordinates on the spaces~$\mathcal H_j$ of Definition~\ref{def:mx-jets} and give each coordinate weight~$j+2$, its scaling degree as a curvature Taylor coefficient.
Substitute the reconstructed potential~\eqref{eq:mx-radial} into a polynomial in the radial gauge fields. This gives the restriction map~$\operatorname{res}_d$ on~$V_d$; its kernel consists exactly of polynomials that vanish on every Maxwell solution.
Write
\begin{equation}\label{eq:mx-polynomials}
\begin{aligned}
 (\operatorname{res}_d v)(Q)&=v\bigl((\kappa_jQ_j)_{j\ge0}\bigr),\\
 \mathcal I_d&=\ker\operatorname{res}_d,\qquad
 \mathcal P_d^0=\operatorname{res}_d(V_d),\qquad
 F_d^0=V_d/\mathcal I_d\simeq\mathcal P_d^0.
\end{aligned}
\end{equation}
For the orthogonal group~$O(4)$ of Definition~\ref{def:uv-actions}, write~$\rho_d^0(h)$ for the action induced on these free polynomial spaces by pullback of the two-form and its Taylor tensors.
Give each~$F_d^0$ its fixed quotient norm from~$V_d$; all vector and operator norms below use this norm on the free factors.
The free unit is the constant class~$\mathbf1=[1]$.
The interacting unit is~$\pi_du_d$ in Definition~\ref{def:iv-fields}; the comparison maps below will send the free unit to that class.
\end{definition}

Wick ordering means subtracting the self-contractions at each insertion. In the resulting correlations, pairings join different insertion points only.
The following finite sum implements this rule.
\begin{definition}[Wick correlations]\label{def:mx-wick}
For monomials~$q_i=\prod_{e\in E_i}p_e$ in the linear probes defined above, regard each repeated factor as a separate occurrence. Let~$E$ be the union of these occurrence sets and let~$i(e)=i$ identify the insertion containing~$e$.
A perfect matching pairs every occurrence with exactly one other occurrence. The disjoint-union convention keeps two copies of the same probe distinct when a monomial repeats it.
Let~$\mathfrak M_{\ne}(E)$ be the matchings that never pair two members of the same~$E_i$.
Set
\begin{equation}\label{eq:mx-wick}
 W_n(q;x)=\sum_{P\in\mathfrak M_{\ne}(E)}
 \prod_{\{e,f\}\in P}K_{p_e,p_f}\bigl(x_{i(e)}-x_{i(f)}\bigr),
 \qquad x\in\Conf n,
\end{equation}
For~$q=(q_1,\ldots,q_n)$, extend this formula complex multilinearly to the polynomial fields of Definition~\ref{def:mx-fields}.
\end{definition}

The configuration space is that of Definition~\ref{def:tests}.
The matching rule treats repeated occurrences separately and is unchanged by their enumeration.
For example, a quadratic insertion has zero one-point expectation, whereas two quadratic insertions give
\begin{equation*}
 W_2(pq,st;x,y)
 =K_{p,s}(x-y)K_{q,t}(x-y)+K_{p,t}(x-y)K_{q,s}(x-y).
\end{equation*}
No value of a kernel at zero is used.
Empty products give~$W_0=1$.\sourcefootnote{\nolinkurl{YM/Maxwell/Equation.lean}: restriction and equation kernel; \nolinkurl{YM/Maxwell/Wick.lean}: matching correlations}

\subsection{Positivity and finite observations}\label{subsec:mx-positive}

This subsection proves positivity for polynomial observations from the analytic result for linear probes.
\begin{definition}[Reflected Maxwell observations]\label{def:mx-reflection}
Fix an axis~$\mu$ and let~$\theta=r_{\mu,0}$ be the reflection of Definition~\ref{def:iv-reflections}.
A positive-side point has~$x^\mu>0$.
Reflect the two-form indices and all Taylor arguments by
\begin{equation*}
 (\theta\cdot Q)_{ab}(z)=\sum_{c,d=1}^4\theta_{ac}\theta_{bd}Q_{cd}(\theta^{-1}z),
 \qquad (R_\theta p)(Q)=p(\theta^{-1}\cdot Q).
\end{equation*}
The colour component is unchanged.
Extend~$R_\theta$ multiplicatively to polynomials; on~$F_d^0$ it is~$\rho_d^0(\theta)$ of Definition~\ref{def:mx-fields}.
For a finite list of probes~$p_\alpha$ at positive-side points~$x_\alpha$, their reflected Gram matrix, the matrix of reflected pairings, is
\begin{equation}\label{eq:mx-linear-gram}
 H_{\alpha\beta}
 =K_{p_\alpha,R_\theta p_\beta}\bigl(x_\alpha-\theta x_\beta\bigr).
\end{equation}
\end{definition}

The reflected point is separated from every positive-side point, even when two of the~$x_\alpha$ coincide.

\begin{lemma}[Reflected linear probes]\label{lem:mx-linear-positive}
The matrix~\eqref{eq:mx-linear-gram} is positive semidefinite.
If all~$x_\alpha=x$ and the probes~$p_\alpha$ are linearly independent as functionals on~$\prod_j\mathcal H_j$, it is positive definite: every nonzero linear combination has strictly positive reflected norm.
\end{lemma}

\gap{The analytic proof of Lemma~\ref{lem:mx-linear-positive} is still missing from this manuscript. It must prove nonnegativity for arbitrary positive-side probe lists and strictness at one point modulo exactly the Maxwell equations. In the strictness argument, the response is the pairing of a linear probe with the curvature field: if it vanishes on the reflected open set, analytic continuation makes it vanish at all separated points, and tests against homogeneous Maxwell solutions must then show that the probe is zero on~$\mathcal H_j$. Positivity alone does not establish this strict assertion.\sourcefootnote{\nolinkurl{YM/Next/MAXWELL/Gram.lean}, lines 52--108; \nolinkurl{YM/Next/MAXWELL/Strict.lean}, especially \texttt{level\_detects}, and its imported energy estimates}}

Once strictness is established, it identifies the quotient~\eqref{eq:mx-polynomials} with the quotient by all separated null correlations, just as required for the interacting fields in~\ref{item:iv-smooth-null}.

An inclusion keeps the same polynomial when the allowed degree is enlarged. A derivative moves the insertion point infinitesimally: the coefficient at degree~$j+1$ contributes to the derivative at degree~$j$, with the factor~$j+1$ fixed by the Taylor factorials.

\begin{definition}[Free field operations]\label{def:mx-operations}
Write~$j^0_{d,e}$ for the inclusion induced by~$V_d\subset V_e$, for~$d\le e$.
For a formal Taylor series~$Q$ as in Definition~\ref{def:mx-jets}, translating the observation point infinitesimally in direction~$\mu$ shifts the Taylor coefficients as follows:
\begin{equation*}
 (\mathcal D_\mu Q)_j=(j+1)Q_{j+1}(e_\mu,\cdot,\ldots,\cdot).
\end{equation*}
Here~$e_\mu$ is the coordinate unit vector and the remaining arguments are the~$j$ symmetric Taylor arguments; form indices are retained.
For a polynomial~$q$, let~$D_jq$ denote its ordinary finite-dimensional derivative in the~$Q_j$ variable and define
\begin{equation}\label{eq:mx-polynomial-derivative}
 (\partial_\mu^0q)(Q)=\sum_j D_jq(Q)\bigl[(\mathcal D_\mu Q)_j\bigr].
\end{equation}
The sum is finite, since~$q$ uses only finitely many~$Q_j$.
\end{definition}

The proof below shows that this operation maps~$\mathcal P_d^0$ into~$\mathcal P_{d+1}^0$ and differentiates its correlations with respect to the insertion point.
In derivative bounds, $D_x^s$ denotes the total Fr\'echet derivative with its operator norm, as in Definition~\ref{def:iv-analysis}.

\begin{proposition}[Free fields]\label{prop:mx-free}
Grant Lemma~\ref{lem:mx-linear-positive}.
The correlations~\eqref{eq:mx-wick} descend faithfully to~$F_d^0$ in~\eqref{eq:mx-polynomials} and have the following properties.
\begin{enumerate}[label=\textup{(\roman*)},ref=\theproposition(\roman*)]
\item\label{item:mx-laws} They are normalised, unital and real analytic on~$\Conf n$; they satisfy the symmetry identities~\eqref{eq:iv-symmetries}, with the full~$O(4)$ action on coefficients, and both reflection-positivity conditions~\eqref{eq:iv-os-positive}.
\item\label{item:mx-calculus} The operations of Definition~\ref{def:mx-operations} give compatible inclusions~$j^0_{d,e}:F_d^0\hookrightarrow F_e^0$ and derivatives~$\partial_\mu^0:F_d^0\to F_{d+1}^0$ satisfying~\eqref{eq:iv-field-derivative} and the conjugation and orientation laws of Section~\ref{subsec:iv-smoothness}.
\item\label{item:mx-bounds} For every fixed field list and~$s\in\N$, there are~$C\ge0$ and~$b\in\N$ with
\begin{equation}\label{eq:mx-bounds}
 \|D_x^sW_n(q;x)\|\le Cr^{-b}
 \quad\bigl(0<r\le1,\ |x_i-x_j|_\infty\ge r\text{ for }i\ne j\bigr).
\end{equation}
\end{enumerate}
\end{proposition}

\begin{proof}
\proofstep{1}{Factor the reflected Gram matrix as a polynomial inner product.}
For a finite collection of products supported on the positive side, use the finite set of all their linear factors as the matrix indices.
Within each product, sum over partial matchings between distinct insertion points, and replace each unpaired occurrence~$e$ by~$\sum_aL_{ae}z_a$.
This produces a polynomial~$Q_\alpha(z)$ for each product~$\alpha$.
Every full matching splits uniquely into a matching within each half and the pairs crossing between them.
Write~$P(z)=\sum_\gamma P_\gamma z^\gamma$, where~$\gamma$ is a multi-index of nonnegative integers and~$z^\gamma=\prod_a z_a^{\gamma_a}$.
The reflected Gram matrix is then
\begin{equation}\label{eq:mx-fock}
 G_{\alpha\beta}
 =\langle Q_\alpha,Q_\beta\rangle_{\rm F},
 \qquad
 \langle P,Q\rangle_{\rm F}
 =\sum_\gamma\gamma!\,\overline{P_\gamma}Q_\gamma,
 \qquad \gamma!=\prod_a\gamma_a!.
\end{equation}
This is the finite-dimensional Gaussian factorisation: the reflected pairing is an ordinary polynomial inner product, and every monomial square has the positive weight~$\gamma!$. Hence the reflected matrix is positive semidefinite.
For a single insertion there are no internal pairs, so~$Q(z)=q(L^{\mathsf T}z)$.
For independent coordinate probes at one point, Lemma~\ref{lem:mx-linear-positive} makes~$L$ invertible and hence
\begin{equation}\label{eq:mx-strict}
 q\ne0\quad\Longrightarrow\quad
 W_2\bigl(R_\theta\bar q,q;\theta x,x\bigr)
 =\sum_\gamma\gamma!\,|Q_\gamma|^2>0.
\end{equation}
Thus the only separated null coefficients in~$V_d$ are~$\mathcal I_d$.
Integrating point Gram matrices against finite sums of positive-side joint tests gives smeared positivity.
The integrals converge absolutely on the separated compact supports, by the explicit kernels.

\proofstep{2}{Read the field laws from the matching formula.}
The constant polynomial and the matching rule give the normalisation and unit identities.
The kernel is invariant under translations and transforms as a two-form under~$O(4)$; permuting occurrences proves permutation symmetry and conjugating coefficients proves reality.
All these identities descend to the quotient.

For differentiation, use the infinitesimal translation of Definition~\ref{def:mx-operations}.
It raises weight by one, obeys the Leibniz rule, and preserves~\eqref{eq:mx-equations}.
The induced derivation preserves the images~$\mathcal P_d^0$, and differentiating a matching sum in slot~$i$ gives
\begin{equation*}
 \partial_{i,\mu}W_n(q;x)=W_n(q_{i\to\partial_\mu^0q_i};x).
\end{equation*}
Inclusions are induced by~$V_d\subset V_e$; faithfulness makes them injective and compatible.
Finally, each derivative of a contraction is homogeneous of a fixed negative degree and bounded on the unit sphere.
The finite sums of products in~\eqref{eq:mx-wick} give~\eqref{eq:mx-bounds}, including the empty-list case.\sourcefootnote{\nolinkurl{YM/Next/MAXWELL/Fock.lean}, \nolinkurl{YM/Next/MAXWELL/Studio.lean} and \nolinkurl{YM/Next/MAXWELL/Core.lean}: polynomial factorisation and strictness; \nolinkurl{YM/Next/MAXWELL/System.lean}, lines 29--114: quotient and derivative operations; \nolinkurl{YM/Next/MAXWELL/Caps.lean}: analyticity and bounds}
\end{proof}

Faithfulness means that every nonzero field is detected by some correlation. At each fixed degree, finitely many observations suffice: they turn convergence of correlations into information about the finite-dimensional coefficient maps.
\begin{definition}[Tested observations]\label{def:mx-tests}
For~$v\in F_d^0$, a fixed list~$b=(b_1,\ldots,b_n)$ of free fields in the other insertion slots, and a separated test~$f$ of Definition~\ref{def:tests}, write
\begin{equation}\label{eq:mx-tests}
 \ell(v)=\int f(x)W_{n+1}(v,b;x)\,dx.
\end{equation}
Here~$(v,b)$ is the list obtained by placing~$v$ first.
\end{definition}

\begin{lemma}[Finite separating tests]\label{lem:mx-tests}
Assume Lemma~\ref{lem:mx-linear-positive}.
For each~$d$, finitely many functionals~\eqref{eq:mx-tests} and a constant~$c_d>0$ satisfy
\begin{equation}\label{eq:mx-detection}
 \max_{1\le a\le s_d}|\ell_a(v)|\ge c_d\|v\|
 \qquad(v\in F_d^0).
\end{equation}
\end{lemma}

\begin{proof}
\proofstep{1}{Reduce the common kernel to zero, then use compactness of the unit sphere.}
Strictness in~\eqref{eq:mx-strict} says every nonzero vector is detected by a functional~\eqref{eq:mx-tests}. Starting with~$F_d^0$, choose one nonzero on the current common kernel at each stage; its dimension drops:
\begin{equation*}
 E_0=F_d^0,\qquad E_{a+1}=E_a\cap\ker\ell_{a+1},
 \qquad E_a\ne0\ \Longrightarrow\ \dim E_{a+1}<\dim E_a.
\end{equation*}
Only finitely many choices are needed.
The maximum in~\eqref{eq:mx-detection} is then a continuous positive function on the compact unit sphere. Its minimum is the positive constant~$c_d$.\sourcefootnote{\nolinkurl{YM/Next/UV_STRESS/Invert.lean}, lines 48--65; \nolinkurl{YM/Work/A04.lean}, lines 14--36}
\end{proof}

\subsection{The comparison estimate}\label{subsec:mx-comparison}

Return to the same limiting correlations~$S$ of Proposition~\ref{prop:iv-smooth} and the normed interacting fields~$F_d$ of Definition~\ref{def:iv-fields}.
A map from a free field to an interacting field is useful only if it compares their correlations.
The needed estimate is uniform in the size of a joint test, including tests at distinct points with equal coordinate projections.
For~$R\ge1$, put
\begin{equation}\label{eq:mx-window}
 K_{R,n}=\{x:|x|\le R,\ |x_i-x_j|_\infty\ge R^{-1}\text{ for }i\ne j\}.
\end{equation}
Every compact subset of~$\Conf n$ lies inside such a set for sufficiently large~$R$.

\begin{definition}[Finite-cutoff comparison]\label{def:mx-comparison}
The comparison data at finite lattice spacing have the following quantifiers.
For every~$D,N_*,R,\varepsilon,\delta$ with~$D,N_*\in\N$, $R\ge1$ and~$\varepsilon,\delta>0$, require a radius~$0<r\le\delta$ and linear maps
\begin{equation}\label{eq:mx-lifts}
 A_{N,p,d}:F_d^0\longrightarrow V_d\qquad(d\le D).
\end{equation}
These coefficient maps choose lattice polynomials representing free fields.
Their operator norms are eventually bounded on the ordered tails of Definition~\ref{def:iv-tails}; they commute with conjugation and signed permutations and preserve the degree-zero unit there.
For tuples with specified degrees, abbreviate~$A_{N,p}q=(A_{N,p,d_i}q_i)_i$.
For a point list~$x$, write~$rx=(rx_1,\ldots,rx_n)$.
For every~$n\le N_*$, $d_i\le D$, fixed tuple~$q_i\in F_{d_i}^0$, separated test~$f$ supported in~$K_{R,n}$, $M>0$ with~$\|f\|_\infty\le M$, and~$\eta>0$, require eventually
\begin{equation}\label{eq:mx-cutoff-comparison}
 \left|C_{N,p}\bigl(A_{N,p}q;f_r\bigr)-\int f(x)W_n(q;x)\,dx\right|
 \le\varepsilon M\prod_i\|q_i\|+\eta,
 \qquad f_r(y)=r^{-4n}f\Bigl(\frac{y}{r}\Bigr).
\end{equation}
The moments and tests are those of Definitions~\ref{def:moments} and~\ref{def:tests}.
The order is: fix accuracy, degree and point-separation window; then choose the radius and coefficient maps; only after that choose any fixed field tuple and test. This order is what makes one comparison map serve every such test.
The cutoff threshold may depend on that tuple, test,~$M$ and~$\eta$; no common cutoff threshold for all tests is required.
\end{definition}

We also need the independent dimension bound
\begin{equation}\label{eq:mx-rank}
 \dim_\C F_d\le\dim_\C F_d^0\qquad(d\in\N).
\end{equation}

\gap{The ultraviolet sections must produce~\eqref{eq:mx-lifts}--\eqref{eq:mx-cutoff-comparison} and~\eqref{eq:mx-rank} for the same Wilson family and the same retained limits as in Section~\ref{sec:infinite-volume}. The comparison includes bounded coefficient maps, symmetries, the degree-zero unit and all the displayed test quantifiers. Centred bounds alone do not provide it.\sourcefootnote{\nolinkurl{YM/Next/UV_STRESS/Station.lean}: exact finite-cutoff comparison condition; \nolinkurl{YM/Next/UV_STRESS/Granular.lean} and \nolinkurl{YM/MergeProof/East/Plume.lean}: suppliers; \nolinkurl{YM/Next/UV_STRESS/Invert.lean}, lines 104--116: the separate dimension input}}

For maps~$U_d:F_d^0\to F_d$, write~$Uq=(U_{d_i}q_i)_i$, just as for the coefficient maps in Definition~\ref{def:mx-comparison}.
First pass that comparison to the already constructed continuum theory.
Bounded coefficient maps have a common cluster limit in finite dimensions.
The extraction is used to estimate the same~$S$, rather than to construct another limiting law. Equal coordinate projections are needed because the comparison tests are supported on configurations where distinct points may share one coordinate; those configurations belong to~$\Conf n$ even though the first compactness extraction used a smaller full-measure set.

\begin{lemma}[Passage to continuum fields]\label{lem:mx-passage}
Assume~\eqref{eq:mx-lifts}--\eqref{eq:mx-cutoff-comparison} and the conclusions of Propositions~\ref{prop:iv-raw}, \ref{prop:iv-smooth} and~\ref{prop:iv-euclidean}.
There are linear maps~$U_d:F_d^0\to F_d$, $d\le D$, commuting with conjugation and signed permutations, with~$U_0\mathbf1=\mathbf1$, such that
\begin{equation}\label{eq:mx-continuum-comparison}
 \left|\int f(x)\bigl[S_n(Uq;rx)-W_n(q;x)\bigr]\,dx\right|
 \le\varepsilon M\prod_i\|q_i\|
\end{equation}
for all the tuples and tests in~\eqref{eq:mx-cutoff-comparison}.
\end{lemma}

\begin{proof}
\proofstep{1}{Extract common finite-dimensional limits and control coefficient replacement.}
More precisely, use a refining ultrafilter as in Definition~\ref{def:iv-tails} and write its limits as~$A_d$.
Expansion in fixed polynomial bases and~\eqref{eq:iv-separated-bound} give, for the fixed rescaled test,
\begin{equation*}
 \left|C_{N,p}(A_{N,p}q;f_r)-C_{N,p}(Aq;f_r)\right|
 \le C_{r,f,q}\sum_{d\le D}\|A_{N,p,d}-A_d\|\longrightarrow0.
\end{equation*}
Put~$U_d=\pi_dA_d$, using the quotient map of Proposition~\ref{prop:iv-smooth}.
The retained ordered limit and the change of variables~$y=rx$ yield
\begin{equation*}
 C_{N,p}(A_{N,p}q;f_r)\longrightarrow\int f(x)S_n(Uq;rx)\,dx.
\end{equation*}
\proofstep{2}{Pass the finite comparison to the retained continuum family.}
All symmetry and unit identities pass through the finite-dimensional limits.
For each fixed test, take the limit in~\eqref{eq:mx-cutoff-comparison}, then let~$\eta\downarrow0$.
The same cluster maps work for every test because they were chosen before any test was fixed. For each fixed tuple and test, the local bound in~\eqref{eq:iv-separated-bound} controls the coefficient replacement; the finite-cutoff inequality then passes to the limit along that already chosen common refinement.\sourcefootnote{\nolinkurl{YM/Next/UV_STRESS/Station.lean}, especially \texttt{zoom\_integral} and \texttt{harvested}}
\end{proof}

\subsection{Compatible field bases}\label{subsec:mx-bases}

We next choose comparisons at successively smaller radii. Compatibility means that if a free polynomial is included from degree~$d$ into degree~$e$, its image is the same interacting field under either map.
The same free field can appear in two degree spaces through an inclusion. The maps of Lemma~\ref{lem:mx-passage} may initially send these two copies to different interacting fields, so we correct the higher-degree map to agree on the already included lower-degree part.


For a radius sequence~$r_k$ and maps~$U_{k,d}:F_d^0\to F_d$, write
\begin{equation}\label{eq:mx-difference}
 E_k(q;x)=S_n(U_kq;r_kx)-W_n(q;x).
\end{equation}
\begin{definition}[Uniform tested convergence]\label{def:mx-tested-convergence}
For every fixed field tuple, the convergence used below is
\begin{equation}\label{eq:mx-weak}
 \sup_{\substack{f\text{ as in Definition }\ref{def:tests}\\
 \supp f\subset K,\ \|f\|_\infty\le1}}
 \left|\int f(x)E_k(q;x)\,dx\right|\longrightarrow0
 \qquad(K\Subset\Conf n),
\end{equation}
Here~$K\Subset\Conf n$ means a compact subset of the separated configuration space.
\end{definition}

The supremum controls every test supported in the same compact region with the same size bound. This uniformity lets us use tests that vary with~$k$ in the derivative argument; convergence for one fixed test would not suffice.

\begin{definition}[Moving tested observations]\label{def:mx-moving-tests}
For a finite list of tested observations~$\ell_a$ on~$F_d^0$ from Definition~\ref{def:mx-tests}, let~$f_a$ be its test and~$b_a$ its list of~$m_a$ other fields.
Given~$r_k$ and maps~$U_{k,d}$, define
\begin{equation}\label{eq:mx-moving-tests}
 \ell_{a,k}(v)=\int f_a(x)S_{m_a+1}\bigl(U_k(v,b_a);r_kx\bigr)\,dx,
 \qquad B_kv=(\ell_{a,k}(v))_a,\quad Bv=(\ell_a(v))_a.
\end{equation}
These are the same measurements, now made on the interacting fields at radius~$r_k$.
\end{definition}

\begin{proposition}[Compatible comparisons]\label{prop:mx-frames}
Assume the finite-cutoff input~\eqref{eq:mx-lifts}--\eqref{eq:mx-cutoff-comparison}, the dimension bound~\eqref{eq:mx-rank}, and the field operations and Euclidean properties of Propositions~\ref{prop:iv-raw}, \ref{prop:iv-smooth}, \ref{prop:iv-euclidean} and~\ref{prop:mx-free}.
There exist~$r_k\downarrow0$ strictly and complex-linear isomorphisms~$U_{k,d}:F_d^0\to F_d$ satisfying~\eqref{eq:mx-weak}.
For each fixed~$d\le e$, eventually in~$k$,
\begin{equation}\label{eq:mx-compatible}
 \begin{aligned}
 U_{k,e}j^0_{d,e}&=j_{d,e}U_{k,d},& U_{k,d}\mathbf1&=\mathbf1,\\
 U_{k,d}\bar q&=\overline{U_{k,d}q},&
 U_{k,d}\rho_d^0(h)&=\rho_d(h)U_{k,d}\quad(h\in H_4).
 \end{aligned}
\end{equation}
\end{proposition}

\begin{proof}
\proofstep{1}{Choose radii that meet successively larger finite requirements.}
Apply Lemma~\ref{lem:mx-passage} with
\begin{equation*}
 D_k=N_{*,k}=k,\qquad R_k=k+1,\qquad
 \varepsilon_k=\frac{1}{k+1},\qquad
 r_{k+1}\le\min\Bigl(\frac{r_k}{2},\frac{1}{k+2}\Bigr).
\end{equation*}
The comparison hypothesis permits the last condition by choosing the next radius bound first.
Each compact separated set eventually fits inside~$K_{R_k,n}$.
For each fixed tuple,~\eqref{eq:mx-continuum-comparison} therefore gives~\eqref{eq:mx-weak} for the initial maps.

\proofstep{2}{Correct each higher-degree map on the included lower-degree space.}
Write~$j_d^0=j^0_{d,d+1}$ and~$j_d=j_{d,d+1}$.
Choose a linear left inverse~$R_d:F_{d+1}^0\to F_d^0$, so~$R_dj_d^0=1$.
It extracts from a higher-degree field the component lying in the copy of degree~$d$. Thus~$1-j_d^0R_d$ removes that old component before the new part is mapped.
It may be chosen real and commuting with the~$H_4$ action of Definition~\ref{def:iv-coefficients}: average any left inverse first with its conjugate and then over~$H_4$.
Both averages remain left inverses because~$j_d^0$ intertwines these symmetries.
For~$d<D_k$, define corrected maps recursively by
\begin{equation}\label{eq:mx-correction}
 \widetilde U_{k,0}=U_{k,0},\qquad
 \widetilde U_{k,d+1}
 =j_d\widetilde U_{k,d}R_d
   +U_{k,d+1}(1-j_d^0R_d).
\end{equation}
Since~$R_dj_d^0=1$, this gives~$\widetilde U_{k,d+1}j_d^0=j_d\widetilde U_{k,d}$.
The symmetry identities pass through the recursion; the degree-zero unit identity propagates to every degree.
For~\eqref{eq:mx-weak}, expand a corrected correlation using
\begin{equation*}
 q=j_d^0R_dq+(q-j_d^0R_dq).
\end{equation*}
After finitely many expansions it is a sum of original comparisons with fixed free fields, possibly in lower degrees.
The inclusion identities in both theories identify the corresponding sum of free correlations with~$W_n(q;x)$.
Thus the corrected maps retain~\eqref{eq:mx-weak}.

\proofstep{3}{Use separating tests and the rank bound to obtain isomorphisms.}
Fix~$d$ and choose the finite functionals of Lemma~\ref{lem:mx-tests}.
Use their maps~$B_k:F_d^0\to\C^{s_d}$ and~$B$ from Definition~\ref{def:mx-moving-tests}, with the corrected comparisons~$\widetilde U_k$.
Then~\eqref{eq:mx-weak}, evaluated on a basis, gives operator-norm convergence to~$B=(\ell_a)_a$.
For large~$k$,~\eqref{eq:mx-detection} implies
\begin{equation*}
 \|B_kv\|_\infty\ge\frac{c_d}{2}\|v\|,
 \qquad
 \widetilde U_{k,d}v=0\ \Longrightarrow\ B_kv=0\ \Longrightarrow\ v=0.
\end{equation*}
Hence~$\widetilde U_{k,d}$ is injective.
Injectivity says the image has at least the free-space dimension. Together with the independent upper bound~\eqref{eq:mx-rank},
\begin{equation*}
 \dim F_d^0\le\dim F_d\le\dim F_d^0,
\end{equation*}
so it is an isomorphism eventually at each fixed degree.
Fill the finitely many exceptional early terms at each degree with any fixed isomorphism; the assertions are eventual.
For a fixed tuple or a fixed pair of degrees these replacements have no effect on the eventual assertions, proving~\eqref{eq:mx-compatible} and~\eqref{eq:mx-weak}.\sourcefootnote{\nolinkurl{YM/Next/UV_STRESS/Rates.lean}: diagonal choice of radii; \nolinkurl{YM/Next/UV_STRESS/Arrange.lean}: equivariant retractions and exact degree correction; \nolinkurl{YM/Next/UV_STRESS/Invert.lean}, especially \texttt{Stations.equivalences}: eventual injectivity and isomorphisms}
\end{proof}

\subsection{Derivatives of the comparison}\label{subsec:mx-smooth}

The tested convergence~\eqref{eq:mx-weak} also determines how the derivatives of the comparison fields change with~$k$.
Integration by parts identifies their values against every separating functional; finite-dimensional invertibility then identifies the derivative maps themselves.
This is the extra argument needed before using a smooth short-distance limit.

\begin{proposition}[Smooth short-distance limits]\label{prop:mx-smooth}
Under the hypotheses of Proposition~\ref{prop:mx-frames}, for every~$d\in\N$ and coordinate direction~$\mu$, the maps there satisfy, in operator norm,
\begin{equation}\label{eq:mx-transported-derivative}
 r_kU_{k,d+1}^{-1}\partial_\mu U_{k,d}\longrightarrow\partial_\mu^0.
\end{equation}
For every fixed tuple~$q$, compact~$K\subset\Conf n$ and~$s\in\N$,
\begin{equation}\label{eq:mx-smooth-limit}
 \sup_{x\in K}\|D_x^s E_k(q;x)\|\longrightarrow0.
\end{equation}
\end{proposition}

\begin{proof}
\proofstep{1}{Recover the transported derivative from finitely many observations.}
Put~$D_{k,d,\mu}=r_kU_{k,d+1}^{-1}\partial_\mu U_{k,d}$. The factor~$r_k$ converts differentiation in the rescaled coordinate~$r_kx$ back to differentiation in the original coordinate.
The field derivative identity~\eqref{eq:iv-field-derivative} and the chain rule give
\begin{equation*}
 \partial_{i,\mu}S_n(U_kq;r_kx)
 =S_n\bigl(U_k(q_{i\to D_{k,d_i,\mu}q_i});r_kx\bigr).
\end{equation*}
Use the finite tested observations at degree~$d+1$ from Lemma~\ref{lem:mx-tests}, with~$\ell_{a,k},f_a,b_a,m_a$ as in Definition~\ref{def:mx-moving-tests}.
Integration by parts in the first insertion variable gives
\begin{equation}\label{eq:mx-tested-derivative}
 \ell_{a,k}(D_{k,d,\mu}q)
 =-\int\partial_{1,\mu}f_a(x)S_{m_a+1}\bigl(U_k(q,b_a);r_kx\bigr)\,dx
 \longrightarrow\ell_a(\partial_\mu^0q).
\end{equation}
The maps~$B_k$ of Definition~\ref{def:mx-moving-tests} converge in operator norm to the injective map~$B=(\ell_a)_a$, as in the preceding proof.
Choose a left inverse~$L$ of~$B$, so~$LB=1$.
Then~$LB_k\to1$, and for large~$k$,
\begin{equation*}
 D_{k,d,\mu}q=(LB_k)^{-1}L\bigl(B_kD_{k,d,\mu}q\bigr)
 \longrightarrow\partial_\mu^0q.
\end{equation*}
Convergence on a finite basis gives~\eqref{eq:mx-transported-derivative}.

\proofstep{2}{Extend tested convergence to every positional derivative.}
Repeated differentiation expresses each positional derivative in terms of higher-degree fields whose free coordinates converge, by~\eqref{eq:mx-transported-derivative}.
Multilinearity and expansion in finite bases extend~\eqref{eq:mx-weak} to these converging coordinate lists.
For a coordinate multi-index~$\alpha$ and compact~$K\Subset\Conf n$, this gives
\begin{equation*}
 \sup_{\substack{f\text{ as in Definition }\ref{def:tests}\\\supp f\subset K,\ \|f\|_\infty\le1}}
 \left|\int f(x)D^\alpha E_k(q;x)\,dx\right|\longrightarrow0.
\end{equation*}
Uniform control of integrals against all tests with sup norm at most one controls the integral of the absolute value: approximate the complex sign of each derivative by smooth tests on a slightly larger compact neighbourhood of~$K$ inside~$\Conf n$. The dual formula for the integral of an absolute value therefore gives
\begin{equation*}
 \|D^\alpha E_k(q;\cdot)\|_{L^1(K)}\longrightarrow0.
\end{equation*}

\proofstep{3}{Use a cutoff and Fourier inversion for uniform derivative bounds.}
Choose a smooth cutoff function~$\chi$ supported in that larger set and equal to one near~$K$. The L1 convergence of every derivative makes the Fourier transform decay uniformly after integration by parts; Fourier inversion then bounds the derivatives pointwise.
Writing~$a=4n$ for the dimension of the point configuration space and~$\alpha$ for a coordinate multi-index, integration by parts in the Fourier transform and Fourier inversion give
\begin{equation*}
 \|D^s(\chi E_k)\|_\infty
 \le C_{\chi,s}\sum_{|\alpha|\le s+a+1}
       \|D^\alpha(\chi E_k)\|_{L^1}\longrightarrow0.
\end{equation*}
This proves~\eqref{eq:mx-smooth-limit}; the zero-dimensional case is just convergence of constants.\sourcefootnote{\nolinkurl{YM/Next/UV_STRESS/Focus.lean}, especially \texttt{head\_tilt}, \texttt{bend\_limit} and \texttt{travels\_jet}; \nolinkurl{YM/Work/A04.lean}, \texttt{tendsto\_of\_fitted}: recovery from varying finite tests}
\end{proof}

A simple nonabelian Lie algebra has no nonzero invariant linear functional: invariance forces~$\lambda([X,Y])=0$, while its commutators span~$\g$.
Apart from constants, each curvature coefficient has weight at least two, so a product of two has weight at least four. Below weight four, only one curvature coefficient or its first derivative can occur; gauge invariance would require an invariant linear functional on~$\g$.

\begin{corollary}[Degrees below four]\label{cor:mx-low}
If~$\g$ is simple in Definition~\ref{def:gauge}, then, for~$d<4$,
\begin{equation}\label{eq:mx-low}
 F_d^0=\C\mathbf1,\qquad \dim_\C F_d^0=1.
\end{equation}
Under~\eqref{eq:mx-rank}, the interacting space also satisfies~$F_d=\C\mathbf1$.
\end{corollary}

\begin{proof}
Simplicity and noncommutativity imply~$[\g,\g]=\g$. Below weight four only a single curvature coefficient or its first derivative could occur; gauge invariance would require a linear functional~$\lambda:\g\to\C$ vanishing on all commutators, and hence~$\lambda=0$:
\begin{equation*}
 \lambda([X,Y])=0\ \text{for all }X,Y\in\g
 \quad\Longrightarrow\quad\lambda=0.
\end{equation*}
Thus~$V_d$ consists of constants, and~$W_1(1)=1$ shows that the constant survives the free quotient.
The dimension bound and the nonzero interacting unit from~\eqref{eq:iv-symmetries} give the last assertion.\sourcefootnote{\nolinkurl{YM/UVNext/FREE_STRESS_CERTIFICATE/Size.lean}, \texttt{just\_one} and \texttt{dim\_before}; \nolinkurl{YM/Trace2/Transfer.lean}, \texttt{first\_grade}}
\end{proof}

At each stage, Proposition~\ref{prop:mx-frames} chooses a radius small enough for finitely many degree and test requirements. The construction does not ensure these radii form a subsequence of a sequence specified in advance.
That stronger choice requires comparison estimates uniform over varying scale intervals, and is reserved for the trace section.
The degree-four observable in Proposition~\ref{prop:trace} is the first nonconstant gauge-invariant field retained by this comparison. Identifying its normalisation requires further identities for differentiated correlations; the radii must also work uniformly across the prescribed scale intervals.
Those later arguments use the field spaces and derivative limits established here.
